% ============================================================
%  Chapter 9 — Stay Safe and Curious
%  The Secret Life of the Internet
% ============================================================

\chapter{Stay Safe and Curious}

\chapteropening{%
  You have learned the secret life of the internet.
  Now it is time to become a smart, safe, and curious explorer!
}
\chapterbadge{Final Quest: Safe Explorer Mode}

% --- Opening story ---
\section*{The Explorer's Badge}

Dot, Bee Browser, Rory Router, Willa Wi-Fi, Sunny Server, Captain Cable,
Scout Search, and Cloud Pal had all gathered together.

They looked at Sam — who had come a very long way since that first question
about the dolphin video.

\character{Dot}{You have followed me through Wi-Fi signals, undersea cables,
  data centres, and packet highways.
  Now it is time for your final challenge:
  learning how to explore the internet \textbf{safely and wisely}.}

Each character reached into a bag and pulled out a tool.
Together they built the \textbf{Digital Explorer's Backpack}.

\sectionrule

% --- Safety tools ---
\section*{The Digital Explorer's Backpack}

\subsection*{Tool 1: A Strong Password}

A \term{password} is like a key to your online account.
A strong password:
\begin{itemize}
  \item Is long (at least 8 characters).
  \item Mixes letters, numbers, and symbols.
  \item Is different for every account.
  \item Is never shared with anyone except a trusted adult.
\end{itemize}

\begin{picturethis}
  Imagine your account is a house.
  A weak password is like leaving the front door open.
  A strong password is a sturdy lock.
  And sharing your password with someone you do not trust is
  like handing them a copy of your key.
\end{picturethis}

\vspace{6pt}

\subsection*{Tool 2: Think Before You Click}

Not every link on the internet is safe.
Some links try to trick you into giving personal information
or downloading something harmful.

\begin{itemize}
  \item If an email or message seems strange, do not click on links in it.
  \item If a website asks for your full name, address, or phone number and you did not expect it — stop.
  \item If something feels wrong, ask a trusted adult.
\end{itemize}

\character{Scout Search}{Good searchers and safe explorers both check twice.
  If something seems too exciting, too scary, or too good to be true —
  it probably needs a second look.}

\vspace{6pt}

\subsection*{Tool 3: The Buddy Rule}

Just as you would not explore a strange place alone in real life,
it is wise to have a trusted adult you can talk to about what you see online.

\begin{itemize}
  \item Tell a parent, carer, or teacher if you see something that worries you.
  \item Ask an adult before signing up for a new website or app.
  \item Remember: you can always close the page or walk away.
\end{itemize}

\vspace{6pt}

\subsection*{Tool 4: Be Kind Online}

The internet connects you to real people all over the world.

\begin{itemize}
  \item Treat people online the same way you would in person — with kindness and respect.
  \item If someone says something unkind to you online, it is not your fault.
        Tell a trusted adult.
  \item Think before posting or sending: \emph{``Would I say this to someone's face?''}
\end{itemize}

\begin{picturethis}
  Words online are real words.
  A kind message can make someone's day.
  An unkind one can really hurt — even when you cannot see the person's face.
\end{picturethis}

\vspace{6pt}

\subsection*{Tool 5: Check Whether Things Are True}

Not everything on the internet is accurate.
Anyone can publish almost anything.

\begin{itemize}
  \item Look for well-known, trusted sources.
  \item If something surprising is claimed, check another reliable source.
  \item Ask a trusted adult or librarian if you are unsure.
\end{itemize}

\vspace{6pt}

\subsection*{Tool 6: Balance Your Screen Time}

The internet is wonderful — but so is the world away from screens.

\begin{itemize}
  \item Take breaks from screens regularly.
  \item Notice how you feel after a long session online.
  \item Balance internet time with physical activity, reading, and face-to-face time with friends and family.
\end{itemize}

\sectionrule

% --- Yes / No choices ---
\section*{Smart Explorer Choices}

\begin{quickmap}
  \textbf{Smart choice or not so smart?}
  \begin{center}
    \begin{tabular}{p{0.7\linewidth}l}
      \textbf{Situation} & \textbf{Smart?} \\
      \hline
      You use a different password for each account. & Yes! \\
      You click a link in an email that says you have won a prize. & Not smart. \\
      You tell a trusted adult when you see something worrying. & Yes! \\
      You share your home address with a stranger online. & Never. \\
      You take a break after an hour of screen time. & Great idea! \\
      You say something unkind in a comment because no one can see you. & Still unkind. \\
      You check a fact in a second source before believing it. & Good thinking! \\
    \end{tabular}
  \end{center}
\end{quickmap}

\sectionrule

% --- Code of Honour ---
\section*{Digital Explorer's Code of Honour}

\begin{picturethis}
  A real explorer does not just know facts — they make promises.
  Read this code aloud and sign it like a true digital adventurer.
\end{picturethis}

\begin{enumerate}[label=\textbf{\arabic*.}]
  \item I will protect my passwords and keep personal information private.
  \item I will pause and think before I click links, downloads, or pop-ups.
  \item I will check facts in more than one trusted place before sharing.
  \item I will be kind in comments, chats, games, and messages.
  \item I will ask a trusted adult for help when something feels wrong.
  \item I will balance screen time with rest, play, and real-world adventures.
\end{enumerate}

\needspace{6\baselineskip}
\vspace{6pt}
\begin{center}
  \begin{tcolorbox}[colback=SunYellow!15,colframe=SunYellow!70!black,arc=8pt,boxrule=1pt,width=0.9\linewidth]
    \textbf{Explorer's Signature:} \rule{0.48\linewidth}{0.4pt}\hfill
    \textbf{Date:} \rule{0.2\linewidth}{0.4pt}
  \end{tcolorbox}
\end{center}

\begin{diagram}[Your six safety tools, all on one explorer's shield.]
  \begin{tikzpicture}[>=Stealth]
    \coordinate (C) at (0,0.2);
    \coordinate (P1) at (0,3.0);
    \coordinate (P2) at (2.1,1.8);
    \coordinate (P3) at (1.8,0.2);
    \coordinate (P4) at (0,-2.8);
    \coordinate (P5) at (-1.8,0.2);
    \coordinate (P6) at (-2.1,1.8);

    \begin{scope}
      \clip (0,3.2) .. controls (2.8,2.8) and (2.4,-0.5) .. (0,-3.2)
                  .. controls (-2.4,-0.5) and (-2.8,2.8) .. (0,3.2) -- cycle;
      \fill[SkyBlue!65] (C) -- (P1) -- (P2) -- cycle;
      \fill[LeafGreen!60] (C) -- (P2) -- (P3) -- cycle;
      \fill[SunYellow!70] (C) -- (P3) -- (P4) -- cycle;
      \fill[WarmOrange!70] (C) -- (P4) -- (P5) -- cycle;
      \fill[SoftPurple!60] (C) -- (P5) -- (P6) -- cycle;
      \fill[CoralRed!55] (C) -- (P6) -- (P1) -- cycle;
    \end{scope}

    \draw[line width=1.2pt,DarkText] (0,3.2) .. controls (2.8,2.8) and (2.4,-0.5) .. (0,-3.2)
         .. controls (-2.4,-0.5) and (-2.8,2.8) .. (0,3.2) -- cycle;
    \draw[DarkText,thick] (C) -- (P1);
    \draw[DarkText,thick] (C) -- (P2);
    \draw[DarkText,thick] (C) -- (P3);
    \draw[DarkText,thick] (C) -- (P4);
    \draw[DarkText,thick] (C) -- (P5);
    \draw[DarkText,thick] (C) -- (P6);

    \node[font=\scriptsize\bfseries,align=center] at (0.8,2.2) {Strong\\Passwords};
    \node[font=\scriptsize\bfseries,align=center] at (1.6,1.0) {Think\\Before\\Clicking};
    \node[font=\scriptsize\bfseries,align=center] at (0.9,-1.0) {Check\\Facts};
    \node[font=\scriptsize\bfseries,align=center] at (-0.9,-1.0) {Screen\\Balance};
    \node[font=\scriptsize\bfseries,align=center] at (-1.6,1.0) {Buddy\\Rule};
    \node[font=\scriptsize\bfseries,align=center] at (-0.8,2.2) {Be Kind\\Online};
  \end{tikzpicture}
\end{diagram}

\sectionrule

% --- Encouragement ---
\section*{You Are a Digital Explorer!}

All the characters gathered around Sam one last time.

\character{Willa Wi-Fi}{You know how I carry signals through the air now!}
\character{Rory Router}{And how I send things the right way!}
\character{Sunny Server}{And how I store all that information!}
\character{Captain Cable}{And how data races under the ocean to reach you!}
\character{Bee Browser}{And how to visit any website you want!}
\character{Cloud Pal}{And what I really am!}
\character{Scout Search}{And how to find the answers you need!}

\character{Dot}{And you know that the internet is not magic —
  it is \textbf{science, engineering, and millions of clever ideas}
  working together to connect the world.
  \textbf{Stay curious.  Stay kind.  Stay safe.}
  And remember — every explorer started by asking one simple question.
  Keep asking yours!}

\sectionrule

\begin{newwords}
  \begin{description}[labelwidth=3.2cm, leftmargin=3.4cm]
    \item[\term{Password}]       A secret code that protects an online account.
    \item[\term{Privacy}]        Keeping personal information to yourself and choosing who sees it.
    \item[\term{Phishing}]       A trick that tries to get you to share personal information by pretending to be trustworthy.
    \item[\term{Screen time}]    The amount of time spent using a screen-based device.
    \item[\term{Digital citizen}] A person who uses the internet responsibly, kindly, and safely.
  \end{description}
\end{newwords}

\sectionrule

\begin{tryit}
  \textbf{Pack your Digital Explorer's Backpack!}

  Draw a backpack and fill it with six items:
  \begin{enumerate}
    \item A padlock — for a \textbf{strong password}.
    \item A magnifying glass — for \textbf{checking facts}.
    \item A shield — for \textbf{thinking before clicking}.
    \item A smiley face — for \textbf{being kind online}.
    \item A pair of hands — for \textbf{the buddy rule} (always have a trusted adult).
    \item A clock — for \textbf{balancing screen time}.
  \end{enumerate}

  \emph{Now your backpack is ready.  Go explore — carefully, kindly, curiously!}
\end{tryit}

\sectionrule

\begin{bigidea}
  The internet is amazing — but \textbf{you} are the most important part.
  Stay \textbf{curious}, stay \textbf{kind}, and stay \textbf{safe}.
  Every great explorer knows the world is full of wonderful things to discover.
\end{bigidea}
